\documentclass[11pt]{article}
\usepackage{amsmath,amssymb}
\title{Room 12 (P2 even-$a$ reciprocal-bimodal mystery): moderator summary, consolidated}
\date{2026-09-27}
\begin{document}
\maketitle

\noindent This room produced several partially-conflicting sub-findings across multiple seat
attempts (some delegated attempts stalled and were retried or completed directly). This summary
reconciles them honestly rather than picking one; the reconciliation itself is the room's main
deliverable.

\section*{What is solid (proved, verified twice independently)}

\textbf{The branch-flip identity is exact.} From $H_o=s^{-1}\sum(\ldots)$, $s=c^{N/2}$: $s$
appears only as the overall prefactor, nowhere inside the sum. So flipping $s\to-s$ gives
$H_o\to-H_o$ exactly, for any $(N,a)$, hence $A_+\leftrightarrow A_-$ and
$A_+/A_-\to(A_+/A_-)^{-1}$ exactly. This was derived independently twice (two seat attempts) and
verified numerically to machine/multi-digit precision (Reviewer AW's own independent Python
reimplementation, and the original seat's modified Rust tool) at several $(N,a)$. \textbf{This part
is PROVED} -- it is also, as both seats and Reviewer AW note, algebraically trivial once the formula
is examined ($x\leftrightarrow1/x$ under any sign flip of a quantity appearing only as a divisor).

\section*{What is NOT solid: whether this explains Room 7's observed clustering}

Reviewer AW found the proposed discriminant (a specific $a\bmod8$ branch-selection rule) fails
outside a narrow boundary region at $q=11$ -- a direct counterexample (same residue class, opposite
cluster, well past small-$N$ effects) -- consistent with the branch selection being a genuine
continuous-in-$(N,a)$ effect (a $\sqrt{\ }$ branch-cut crossing) rather than a fixed arithmetic law.
A separate, independent numerical push (seat Ramanujan, 46+ new points per seed across $q=9,11,15,
21,25,29,33$) complicates the picture further and in places directly conflicts with the original
Room 7 description:

\begin{itemize}
\item \textbf{$q=11$: genuinely different from Room 7's ``possibly under-sampled reciprocal-bimodal''
guess.} Even-$a$ is chaotic for small/moderate $N$ (spanning $0.04$ to $117$) but settles by
$N\gtrsim5\times10^5$ into single-valued convergence near $25.4$--$25.7$ -- \emph{no} reciprocal
$\approx1/25$ branch appears at large $N$. A fourth qualitative regime, not bimodal.
\item \textbf{$q=9$: the previously-reported $\approx0.15$ reciprocal sub-cluster did not reproduce}
in 30+ new even-$a$ points (all landed in $6.40$--$7.01$). Flagged as an unresolved discrepancy
with Room 7's original report, not yet explained (possibly Room 7's original $0.15$ points were
themselves few and at small $N$, an under-sampling in the other direction).
\item \textbf{Exactness of the reciprocal PAIRS (as opposed to the branch-flip identity) is only
approximate}, not exact: cross-products of the $q=29$ cluster values (its cleanest case) deviate
from $1$ by $0.3$--$1.3\%$, far short of the precision available. This is expected and NOT a
contradiction of the proved identity above: the identity is about the SAME $(N,a)$ under a sign
flip, not about two DIFFERENT sampled $(N,a)$ points that happen to land near the two branches --
those are different mathematical objects, only approximately reciprocal because they approximately
sit near the same $|H_e/H_o|$ ratio.
\item \textbf{No universality:} $q=33$ shows clean single-branch convergence (no split); $q=21$ is
erratic on both parities with no split; $q=15$ shows a distinct, precision-verified (200 vs.\
1600-bit, bit-identical) near-singular blow-up spanning 100+ orders of magnitude, tied to having
two small prime factors ($3,5$), unrelated to reciprocal-bimodality.
\end{itemize}

A separate seat found a related, complementary formula ($A_+/A_-=-i\cot(\varphi/2)$ when
$|H_e|\approx|H_o|$, $\varphi=\arg H_e-\arg H_o$) that reproduces measured values quantitatively at
some points and observed $\arg H_o$ jumping by $\approx\pi$ every second $a$ -- but this was
observed on a narrower sample and has not been reconciled against Ramanujan's broader, partly
contradictory survey above; it should not be treated as an established law without further checking
against the wider $q$ range.

\section*{Room verdict: ADVANCED on the identity, but the phenomenon is genuinely more complex and
less universal than Room 7 first suggested}

\textbf{What to keep:} the branch-flip identity (proved, real, if algebraically modest).
\textbf{What to revise:} Room 7's ``reciprocal-bimodal at $q=9,25,29$, structure at $q=11$ likely
under-sampled'' summary should be read as \emph{provisional and seed-dependent}, not a settled
pattern -- the broader survey here (Ramanujan) found $q=9$'s specific claimed cluster did not
reproduce, $q=11$ resolves to plain convergence at large $N$, and three further seeds
($15,21,33$) each show yet another, distinct qualitative behavior. \textbf{No general theorem, no
universal mechanism, and no closed form was established.} The honest state of P2's amplitude
question after 3 rooms (3, 7, 12) is: richer and less uniform empirically than any single room's
snapshot suggested, with one small proved algebraic fact (the branch-flip identity) and no
resolved discriminant for which regime a given seed falls into. Repo-only; not paper-citable.

\section*{Process note}

Several delegated sub-agent attempts on this exact topic returned only ``launched in background,
will report later'' with no concrete output before eventually completing (or being redone
directly) -- a recurring failure mode this session, now confirmed multiple times. Future work on
this thread should expect this and verify agent outputs contain real, checkable numbers before
treating any report as final, exactly as was done here (which is why the room's overall picture is
more cautious than any single seat's own conclusion).

\end{document}
